\section{Daftar (list) dan Tabel: Styling Dasar}

Daftar (\code{<ul>} untuk unordered, \code{<ol>} untuk ordered) dan tabel (\code{<table>}) adalah elemen HTML yang sering membutuhkan styling khusus agar sesuai dengan desain halaman. CSS menyediakan properti-properti yang mengontrol tampilan bullet/nomor pada daftar, serta border, spacing, dan background pada tabel. Styling yang tepat menjadikan konten lebih terbaca dan menarik \cite{mdnCss}.

\subsection{Styling Daftar}

Properti \code{list-style-type} mengontrol jenis penanda: \code{disc} (lingkaran solid, default untuk \code{ul}), \code{circle}, \code{square}, \code{decimal}, \code{lower-alpha}, \code{none} (tanpa penanda). Properti \code{list-style-position} menentukan apakah penanda di luar (\code{outside}, default) atau di dalam (\code{inside}) area konten. Properti \code{list-style-image} mengganti penanda dengan gambar kustom. Menghilangkan penanda dengan \code{list-style: none} adalah teknik dasar untuk membuat menu navigasi dari daftar \cite{webdevCss}.

\subsection{Styling Tabel}

Properti \code{border-collapse} mengontrol bagaimana border sel berinteraksi: \code{collapse} menggabungkan border yang bersebelahan menjadi satu, \code{separate} (default) mempertahankan border terpisah dengan jarak yang diatur \code{border-spacing}. Baris bergantian (\textit{zebra striping}) dibuat dengan \code{tr:nth-child(even)} atau \code{tr:nth-child(odd)} untuk meningkatkan keterbacaan tabel yang panjang. Properti \code{text-align} dan \code{vertical-align} mengatur perataan teks di dalam sel \cite{cssTricks}.

\begin{contoh}
\textbf{Studi Kasus: Menu Navigasi Horizontal dari List}

Menu navigasi web umumnya dibangun dari elemen \code{<ul>} yang di-styling dengan CSS. Langkah-langkah:
\begin{enumerate}
  \item Hilangkan bullet dan padding default dengan \code{list-style: none; padding: 0;}
  \item Ubah \code{<li>} menjadi \code{display: inline-block} agar berjajar horizontal
  \item Styling \code{<a>} di dalam menu: hapus underline, tambahkan padding untuk area klik yang besar, dan atur warna hover
  \item Tambahkan background dan border-radius untuk tampilan tombol
\end{enumerate}
\end{contoh}

% Webcode diganti dengan deskripsi ringkas
\textbf{Contoh CSS:} \code{.nav-menu} dengan \code{list-style: none} dan \code{display: flex}; \code{tr:nth-child(even)} untuk zebra striping.

\begin{catatan}
\textbf{Aksesibilitas Menu Navigasi:} Saat mengubah \code{<ul>} menjadi menu navigasi visual, pastikan elemen semantik tetap utuh. Bungkus menu dengan \code{<nav>} dan beri \code{aria-label}. Hindari menghilangkan \code{<li>} atau mengganti \code{<ul>} dengan \code{<div>}---pembaca layar mengandalkan daftar untuk mengumumkan jumlah item menu. Fokus keyboard harus berfungsi pada setiap tautan \cite{mdnAccessibility}.
\end{catatan}

